\documentclass[exercice]{thedocument}%
\usepackage{texmaths-tikz}%
\usepackage{tabularray}%
\nodata
\begin{document}%
    \>S1=A;\>S2=B;\>S3=C;%
    \>L1=3 cm;\>L2=4 cm;\>L3=5 cm;%
    \>S4=D;\>S5=E;\>S6=F;%
    \>L4=20 cm;\>L5=2,1 dm;\>L6=2,9 dm;%
    \>S7=D;\>S8=E;\>S9=F;%
    \>L7=5 m;\>L8=12 m;\>L9=130 dm;%
    \>S10=D;\>S11=E;\>S12=F;%
    \>L10=2,5 km;\>L11=6 km;\>L12=6,5 km;%
    \enonce{0129}%
\end{document}